\mode*
\section{The file system}
\label{sec:filesystem}

The second educational objective is the hierarchical file system: files and
directories, the tree, the working directory, and absolute versus relative
paths.

\subsection{What the literature supports}

The personal-information-management literature shows that fluent users
retrieve files overwhelmingly by \emph{navigating} folder hierarchies
(56--68\% of retrievals) and search only as a last resort (4--15\%), a
preference that better search engines did not change
\autocite{Bergman2008Navigation}.
Folder navigation recruits the brain's spatial-navigation structures, while
query-based search recruits linguistic processing
\autocite{Benn2015BrainFolders}: the file tree is effectively a
\emph{spatial} skill.
Folders also serve thinking and organising, not just retrieval; users
forced to work folder-free could not replace them with search and recency
lists \autocite{Whitham2017Folder}, and location-based finding is preferred
for its reminding function \autocite{BarreauNardi1995}.
The mismatch runs deeper than retrieval: users hold rich mental models of
their files---projects, documents evolving over time---which the operating
system's named-collection-of-bytes view simply ignores
\autocite{May2019FileLifecycle}.

The popular claim that students who grew up with search no longer
understand files and folders is, so far, journalistic rather than
peer-reviewed \autocite[flagged as anecdote]{Chin2021FileNotFound}.
What is backed is weaker but sufficient for our purposes: the file tree is
a spatial skill that the shell \emph{assumes}, and that our teaching must
therefore build.
% TODO(#3): Analyse the course's own grading data (path errors in the terminal
% lab) to substantiate---or refute---the local version of the cohort
% claim; cf. introtools issue #50.

\mode<presentation>{%
\begin{frame}
  \begin{block}{The file tree is a spatial skill}
    \begin{itemize}
      \item Fluent users navigate (56--68\%), search is last resort
        (4--15\%) \autocite{Bergman2008Navigation}.
      \item Navigation uses spatial brain structures
        \autocite{Benn2015BrainFolders}.
      \item Folders support thinking, not just retrieval
        \autocite{Whitham2017Folder}.
      \item \enquote{Students can't use folders}: journalism, not yet
        research \autocite{Chin2021FileNotFound}.
    \end{itemize}
  \end{block}
\end{frame}%
}

\subsection{Candidate critical aspects}
\label{sec:filesystem-aspects}

The literature above describes what fluent users \emph{have}---a spatial
model of a tree they navigate---rather than what learners get wrong.
The aspects below are therefore read off the difference between that
model and the learner who lacks it, with the course's own path errors
(\cref{sec:method-vt-analysis}) as the local evidence that the lack is
real.

\begin{description}
  \item[Everything has a place]
    Every file lives at exactly one path in one tree, whichever program
    made it.
    Read off the spatial model of the fluent user
    \autocite{Bergman2008Navigation,Benn2015BrainFolders} and its
    documented mismatch with how users otherwise think of their
    files---as projects and documents rather than as named bytes at a
    path \autocite{May2019FileLifecycle}: a learner whose files live
    \enquote{in Word} or \enquote{in Downloads, somewhere} has not
    discerned that each has a place in the tree.
    Tested by the item \emph{where a saved file is}
    (\cref{app:knowledge-items}).
  \item[Where am I?]
    Every shell has a current working directory, and commands act
    relative to it.
    The notion has no counterpart in the graphical file manager and is
    not in the literature; it is read off the course's path
    errors---commands run from the wrong directory---and the object of
    learning itself.
    Tested by the item \emph{which directory the command uses}.
  \item[One tree, two views]
    The file manager and \texttt{cd}/\texttt{ls} traverse the same tree.
    The file-system instance of the command model's \emph{same system as
    the GUI} (\cref{sec:shell-aspects}), sharing its status and its item
    \emph{the same files from two sides}.
  \item[Absolute versus relative]
    The same file has many names, depending on where you stand.
    Read off the same path errors---the relative path that was right
    from one directory and is wrong from another---and the consequence
    of the two aspects before it.
    Tested by the item \emph{two names for one file}.
\end{description}

\subsection{Patterns of variation}

\begin{itemize}
  \item Contrast (program varies, place invariant): save a file from the
    editor, download one in the browser, \texttt{touch} one in the
    shell---then \texttt{ls} shows all three at their paths in the one
    tree.
  \item Contrast (view varies, tree invariant): open a directory in the
    file manager and \texttt{ls} it in the shell; create a file in one
    view, observe it in the other.
  \item Contrast (standpoint varies, target invariant): reach the same
    file from different working directories---the relative path varies,
    the absolute path does not.
  \item Fusion: a task that requires moving, listing and referring across
    directories at once.
\end{itemize}
